\newpage
\begin{center}
  \textbf{\large ЗАКЛЮЧЕНИЕ}
\end{center}
\refstepcounter{chapter}
\addcontentsline{toc}{chapter}{ЗАКЛЮЧЕНИЕ}

Выполненный теоретический анализ позволяет сделать несколько итоговых выводов. Цифровизация предприятия общественного питания должна рассматриваться не как совокупность разрозненных программных модулей, а как многослойная система, в которой операционный, учетный, клиентский, интеграционно-аналитический и правовой контуры взаимосвязаны и предъявляют различные требования к проектированию информационной системы.

Проведенное исследование показывает, что типовые рыночные решения для общественного питания в наибольшей степени ориентированы на транзакционное обслуживание, учет и повторные коммуникации с гостем. Вместе с тем для ООО <<\_\_\_\_>> ключевое значение приобретает информационная система в форме интерактивного веб-приложения, которая должна поддерживать путь гостя до визита, во время посещения и после него, обеспечивать накопление значимых клиентских данных и сохранять связь цифрового сервиса с тематическим опытом заведения.

На основе рассмотренных теоретических подходов обосновано, что проектирование такой системы требует объединения нескольких оснований: сервисного дизайна, моделей клиентского пути, представлений о клиентской вовлеченности, требований к работе с данными и правовых ограничений обработки персональных данных. Следовательно, клиентский цифровой контур тематического кафе выступает проектным слоем разрабатываемой информационной системы и не сводится ни к электронному меню, ни к стандартной программе лояльности, ни к кассовой автоматизации.

Таким образом, результаты проведенного теоретического исследования формируют методическую основу анализа предметной области, конкретизации функциональных и нефункциональных требований, выбора архитектурных решений, проектирования структуры данных и пользовательских сценариев будущей системы.
